\documentclass{amsart}
% For dealing with references we use the comment environment
\usepackage{verbatim}
\newenvironment{reference}{\comment}{\endcomment}
%\newenvironment{reference}{}{}
\newenvironment{slogan}{\comment}{\endcomment}
\newenvironment{history}{\comment}{\endcomment}

% For commutative diagrams we use Xy-pic
\usepackage[all]{xy}

% We use 2cell for 2-commutative diagrams.
\xyoption{2cell}
\UseAllTwocells

% We use multicol for the list of chapters between chapters
\usepackage{multicol}

% This is generally recommended for better output
\usepackage{lmodern}
\usepackage[T1]{fontenc}

% For cross-file-references

% Package for hypertext links:
\usepackage{hyperref}

% For any local file, say "hello.tex" you want to link to please

% Theorem environments.
%
\theoremstyle{plain}
\newtheorem{theorem}[subsection]{Theorem}
\newtheorem{proposition}[subsection]{Proposition}
\newtheorem{lemma}[subsection]{Lemma}

\theoremstyle{definition}
\newtheorem{definition}[subsection]{Definition}
\newtheorem{example}[subsection]{Example}
\newtheorem{exercise}[subsection]{Exercise}
\newtheorem{situation}[subsection]{Situation}

\theoremstyle{remark}
\newtheorem{remark}[subsection]{Remark}
\newtheorem{remarks}[subsection]{Remarks}

\numberwithin{equation}{subsection}

% Macros
%
\def\lim{\mathop{\mathrm{lim}}\nolimits}
\def\colim{\mathop{\mathrm{colim}}\nolimits}
\def\Spec{\mathop{\mathrm{Spec}}}
\def\Hom{\mathop{\mathrm{Hom}}\nolimits}
\def\Ext{\mathop{\mathrm{Ext}}\nolimits}
\def\SheafHom{\mathop{\mathcal{H}\!\mathit{om}}\nolimits}
\def\SheafExt{\mathop{\mathcal{E}\!\mathit{xt}}\nolimits}
\def\Sch{\mathit{Sch}}
\def\Mor{\mathop{\mathrm{Mor}}\nolimits}
\def\Ob{\mathop{\mathrm{Ob}}\nolimits}
\def\Sh{\mathop{\mathit{Sh}}\nolimits}
\def\NL{\mathop{N\!L}\nolimits}
\def\CH{\mathop{\mathrm{CH}}\nolimits}
\def\proetale{{pro\text{-}\acute{e}tale}}
\def\etale{{\acute{e}tale}}
\def\QCoh{\mathit{QCoh}}
\def\Ker{\mathop{\mathrm{Ker}}}
\def\Im{\mathop{\mathrm{Im}}}
\def\Coker{\mathop{\mathrm{Coker}}}
\def\Coim{\mathop{\mathrm{Coim}}}

% Boxtimes
%
\DeclareMathSymbol{\boxtimes}{\mathbin}{AMSa}{"02}

%
% Macros for moduli stacks/spaces
%
\def\QCohstack{\mathcal{QC}\!\mathit{oh}}
\def\Cohstack{\mathcal{C}\!\mathit{oh}}
\def\Spacesstack{\mathcal{S}\!\mathit{paces}}
\def\Quotfunctor{\mathrm{Quot}}
\def\Hilbfunctor{\mathrm{Hilb}}
\def\Curvesstack{\mathcal{C}\!\mathit{urves}}
\def\Polarizedstack{\mathcal{P}\!\mathit{olarized}}
\def\Complexesstack{\mathcal{C}\!\mathit{omplexes}}
% \Pic is the operator that assigns to X its picard group, usage \Pic(X)
% \Picardstack_{X/B} denotes the Picard stack of X over B
% \Picardfunctor_{X/B} denotes the Picard functor of X over B
\def\Pic{\mathop{\mathrm{Pic}}\nolimits}
\def\Picardstack{\mathcal{P}\!\mathit{ic}}
\def\Picardfunctor{\mathrm{Pic}}
\def\Deformationcategory{\mathcal{D}\!\mathit{ef}}

\setlength{\emergencystretch}{3em}
\begin{document}
\title[Stacks Project, Tag 07E5]{1744 --- For Noetherian local algebras over characteristic-p fields, geometric regularity is equivalent to regularity and canonical cotangent or absolute-differential injectivity}
\maketitle
\noindent Copyright (C) 2005--2025 Johan de Jong. Stacks Project authors. Source: \href{https://stacks.math.columbia.edu/tag/07E5}{Tag 07E5}.

\noindent Editorial difficulty note (not author text). Difficulty H2: The author characterizes geometric regularity of an arbitrary Noetherian local algebra over an imperfect characteristic-p field through cotangent H1 and absolute differential injectivity. Its hard direction uses independent p-root field extensions, explicit zero-differential cotangent presentations, Cartier dimension equalities and two coupled conormal diagrams to recover differential independence; the opposite direction controls embedding dimension under every finite purely inseparable extension. These inseparability/cotangent dimension arguments are substantially beyond ordinary Jacobian checks..

\noindent Editorial provenance note (not author text). History: excerpt from revision \texttt{a0444\allowbreak 6e57e\allowbreak c1fbc\allowbreak 252a8\allowbreak 71afc\allowbreak ec775\allowbreak 2fb28\allowbreak 07b14}, more-algebra.tex, lines 8696--8848. Reference presentation was adapted; original-author context and core support blocks were retained or added. The mathematical wording is unchanged. Licensed under GNU FDL 1.2 or later; no invariant sections or cover texts. See ../licenses/COPYING. Context blocks reproduce author definitions and supporting results; they are not additional counted problems.

\begin{definition}
\label{algebra-definition-geometrically-regular}
Let $k$ be a field. Let $R$ be a Noetherian $k$-algebra.
The $k$-algebra $R$ is called {\it geometrically regular} over $k$ if
the equivalent conditions of Lemma \ref{algebra-lemma-geometrically-regular} hold.
\end{definition}

\begin{definition}
\label{algebra-definition-regular-local}
Let $(R, \mathfrak m)$ be a Noetherian local ring of dimension $d$.
\begin{enumerate}
\item A {\it system of parameters of $R$} is a sequence of elements
$x_1, \ldots, x_d \in \mathfrak m$ which generates an ideal of
definition of $R$,
\item if there exist $x_1, \ldots, x_d \in \mathfrak m$
such that $\mathfrak m = (x_1, \ldots, x_d)$ then we call
$R$ a {\it regular local ring} and $x_1, \ldots, x_d$ a {\it regular
system of parameters}.
\end{enumerate}
\end{definition}

\begin{definition}
\label{definition-local-complete-intersection}
A ring map $A \to B$ is called a {\it local complete intersection}
if it is of finite type and for some (equivalently any) presentation
$B = A[x_1, \ldots, x_n]/I$ the ideal $I$ is Koszul-regular.
\end{definition}

\begin{definition}
\label{definition-koszul-regular-sequence}
Let $R$ be a ring. Let $r \geq 0$ and let $f_1, \ldots, f_r \in R$
be a sequence of elements. Let $M$ be an $R$-module.
The sequence $f_1, \ldots, f_r$ is called
\begin{enumerate}
\item {\it $M$-Koszul-regular} if
$H_i(K_\bullet(f_1, \ldots, f_r) \otimes_R M) = 0$ for
all $i \not = 0$,
\item {\it $M$-$H_1$-regular} if
$H_1(K_\bullet(f_1, \ldots, f_r) \otimes_R M) = 0$,
\item {\it Koszul-regular} if $H_i(K_\bullet(f_1, \ldots, f_r)) = 0$ for
all $i \not = 0$, and
\item {\it $H_1$-regular} if $H_1(K_\bullet(f_1, \ldots, f_r)) = 0$.
\end{enumerate}
\end{definition}

\begin{definition}
\label{definition-regular-ideal}
Let $R$ be a ring and let $I \subset R$ be an ideal.
\begin{enumerate}
\item We say $I$ is a {\it regular ideal} if for every
$\mathfrak p \in V(I)$ there exists a $g \in R$, $g \not \in \mathfrak p$
and a regular sequence $f_1, \ldots, f_r \in R_g$ such that $I_g$
is generated by $f_1, \ldots, f_r$.
\item We say $I$ is a {\it Koszul-regular ideal} if for every
$\mathfrak p \in V(I)$ there exists a $g \in R$, $g \not \in \mathfrak p$
and a Koszul-regular sequence $f_1, \ldots, f_r \in R_g$ such that $I_g$
is generated by $f_1, \ldots, f_r$.
\item We say $I$ is a {\it $H_1$-regular ideal} if for every
$\mathfrak p \in V(I)$ there exists a $g \in R$, $g \not \in \mathfrak p$
and an $H_1$-regular sequence $f_1, \ldots, f_r \in R_g$ such that $I_g$
is generated by $f_1, \ldots, f_r$.
\item We say $I$ is a {\it quasi-regular ideal} if for every
$\mathfrak p \in V(I)$ there exists a $g \in R$, $g \not \in \mathfrak p$
and a quasi-regular sequence $f_1, \ldots, f_r \in R_g$ such that $I_g$
is generated by $f_1, \ldots, f_r$.
\end{enumerate}
\end{definition}

\begin{definition}
\label{algebra-definition-naive-cotangent-complex}
Let $R \to S$ be a ring map. The {\it naive cotangent complex}
$\NL_{S/R}$ is the chain complex (\ref{algebra-equation-naive-cotangent-complex})
$$
\NL_{S/R} = \left(I/I^2 \longrightarrow \Omega_{R[S]/R} \otimes_{R[S]} S\right)
$$
with $I/I^2$ placed in (homological) degree $1$ and
$\Omega_{R[S]/R} \otimes_{R[S]} S$ placed in degree $0$. We will denote
$H_1(L_{S/R}) = H_1(\NL_{S/R})$\footnote{This module is sometimes
denoted $\Gamma_{S/R}$ in the literature.} the homology in degree $1$.
\end{definition}

\begin{definition}
\label{algebra-definition-regular}
A Noetherian ring $R$ is said to be {\it regular}
if all the localizations $R_{\mathfrak p}$ at primes are
regular local rings.
\end{definition}

\begin{definition}
\label{algebra-definition-separable-field-extension}
Let $K/k$ be a field extension.
\begin{enumerate}
\item We say $K$ is {\it separably generated over $k$} if there exists
a transcendence basis $\{x_i; i \in I\}$ of $K/k$ such that the extension
$K/k(x_i; i \in I)$ is a separable algebraic extension.
\item We say $K$ is {\it separable over $k$} if for every subextension
$k \subset K' \subset K$ with $K'$ finitely generated
over $k$, the extension $K'/k$ is separably generated.
\end{enumerate}
\end{definition}

\begin{definition}
\label{algebra-definition-perfect}
Let $k$ be a field. We say $k$ is {\it perfect}
if every field extension of $k$ is separable over $k$.
\end{definition}

\begin{lemma}[Cartier equality]
\label{lemma-cartier-equality}
Let $K/k$ be a finitely generated field extension.
Then $\Omega_{K/k}$ and $H_1(L_{K/k})$ are finite dimensional and
$\text{trdeg}_k(K) = \dim_K \Omega_{K/k} - \dim_K H_1(L_{K/k})$.
\end{lemma}

\begin{proof}
We can find a global complete intersection
$A = k[x_1, \ldots, x_n]/(f_1, \ldots, f_c)$
over $k$ such that
$K$ is isomorphic to the fraction field of $A$, see
Algebra, Lemma \ref{algebra-lemma-colimit-syntomic}
and its proof. In this case we see that $\NL_{K/k}$ is homotopy
equivalent to the complex
$$
\bigoplus\nolimits_{j = 1, \ldots, c} K \longrightarrow
\bigoplus\nolimits_{i = 1, \ldots, n} K\text{d}x_i
$$
by Algebra, Lemmas \href{https://stacks.math.columbia.edu/tag/00S1}{lemma NL homotopy (Tag 00S1)} and
\ref{algebra-lemma-localize-NL}.
The transcendence degree of $K$ over $k$ is the dimension of $A$
(by Algebra, Lemma \href{https://stacks.math.columbia.edu/tag/00P0}{lemma dimension prime polynomial ring (Tag 00P0)})
which is $n - c$ and we win.
\end{proof}


\begin{lemma}
\label{lemma-transitivity-gamma}
Let $M/L/K$ be field extensions. Then the Jacobi-Zariski
sequence
$$
0 \to H_1(L_{L/K}) \otimes_L M \to
H_1(L_{M/K}) \to
H_1(L_{M/L}) \to
\Omega_{L/K} \otimes_L M \to
\Omega_{M/K} \to
\Omega_{M/L} \to 0
$$
is exact.
\end{lemma}

\begin{proof}
Combine
Lemma \ref{lemma-transitive-colimit-lci-at-end}
with
Algebra, Lemma \ref{algebra-lemma-colimit-syntomic}.
\end{proof}


\begin{lemma}
\label{lemma-gamma-commutative-diagram}
Given a commutative diagram of fields
$$
\xymatrix{
K \ar[r] & K' \\
k \ar[u] \ar[r] & k' \ar[u]
}
$$
with $k'/k$ and $K'/K$ finitely generated field extensions
the kernel and cokernel of the maps
$$
\alpha : \Omega_{K/k} \otimes_K K' \to \Omega_{K'/k'}
\quad\text{and}\quad
\beta : H_1(L_{K/k}) \otimes_K K' \to H_1(L_{K'/k'})
$$
are finite dimensional and
$$
\dim \Ker(\alpha) - \dim \Coker(\alpha)
-\dim \Ker(\beta) + \dim \Coker(\beta)
=
\text{trdeg}_k(k') - \text{trdeg}_K(K')
$$
\end{lemma}

\begin{proof}
The Jacobi-Zariski sequences for $k \subset k' \subset K'$ and
$k \subset K \subset K'$ are
$$
0 \to H_1(L_{k'/k}) \otimes K'  \to
H_1(L_{K'/k}) \to
H_1(L_{K'/k'}) \to
\Omega_{k'/k} \otimes K' \to
\Omega_{K'/k} \to
\Omega_{K'/k'} \to 0
$$
and
$$
0 \to H_1(L_{K/k}) \otimes K' \to
H_1(L_{K'/k}) \to
H_1(L_{K'/K}) \to
\Omega_{K/k} \otimes K' \to
\Omega_{K'/k} \to
\Omega_{K'/K} \to 0
$$
By
Lemma \ref{lemma-cartier-equality}
the vector spaces $\Omega_{k'/k}$, $\Omega_{K'/K}$, $H_1(L_{K'/K})$, and
$H_1(L_{k'/k})$ are finite dimensional and the alternating sum of their
dimensions is $\text{trdeg}_k(k') - \text{trdeg}_K(K')$.
The lemma follows.
\end{proof}


\begin{lemma}
\label{algebra-lemma-colimit-syntomic}
Let $K/k$ be a field extension.
Then $K$ is a filtered colimit of global complete intersection
algebras over $k$. If $K/k$ is separable, then $K$ is a filtered
colimit of smooth algebras over $k$.
\end{lemma}

\begin{proof}
Suppose that $E \subset K$ is a finite subset. It suffices to show that
there exists a $k$ subalgebra $A \subset K$ which contains $E$
and which is a global complete intersection (resp.\ smooth) over $k$.
The separable/smooth case follows from
Lemma \href{https://stacks.math.columbia.edu/tag/037X}{lemma localization smooth separable (Tag 037X)}.
In general let $L \subset K$ be the subfield generated by $E$.
Pick a transcendence basis $x_1, \ldots, x_d \in L$ over $k$.
The extension $L/k(x_1, \ldots, x_d)$ is finite.
Say $L = k(x_1, \ldots, x_d)[y_1, \ldots, y_r]$.
Pick inductively polynomials $P_i \in k(x_1, \ldots, x_d)[Y_1, \ldots, Y_r]$
such that $P_i = P_i(Y_1, \ldots, Y_i)$ is monic in $Y_i$ over
$k(x_1, \ldots, x_d)[Y_1, \ldots, Y_{i - 1}]$ and maps to the
minimum polynomial of $y_i$ in
$k(x_1, \ldots, x_d)[y_1, \ldots, y_{i - 1}][Y_i]$.
Then it is clear that $P_1, \ldots, P_r$ is a regular sequence
in $k(x_1, \ldots, x_r)[Y_1, \ldots, Y_r]$ and that
$L = k(x_1, \ldots, x_r)[Y_1, \ldots, Y_r]/(P_1, \ldots, P_r)$.
If $h \in k[x_1, \ldots, x_d]$ is a polynomial such that
$P_i \in k[x_1, \ldots, x_d, 1/h, Y_1, \ldots, Y_r]$, then
we see that $P_1, \ldots, P_r$ is a regular sequence in
$k[x_1, \ldots, x_d, 1/h, Y_1, \ldots, Y_r]$ and
$A = k[x_1, \ldots, x_d, 1/h, Y_1, \ldots, Y_r]/(P_1, \ldots, P_r)$
is a global complete intersection. After adjusting our choice of $h$
we may assume $E \subset A$ and we win.
\end{proof}


\begin{lemma}
\label{algebra-lemma-size-extension-pth-roots}
Let $k$ be a field of characteristic $p > 0$.
Let $a_1, \ldots, a_n \in k$ be elements such that
$\text{d}a_1, \ldots, \text{d}a_n$ are linearly independent in
$\Omega_{k/\mathbf{F}_p}$. Then the field extension
$k(a_1^{1/p}, \ldots, a_n^{1/p})$ has degree $p^n$ over $k$.
\end{lemma}

\begin{proof}
By induction on $n$. If $n = 1$ the result is
Lemma \ref{algebra-lemma-derivative-zero-pth-power}.
For the induction step, suppose that $k(a_1^{1/p}, \ldots, a_{n - 1}^{1/p})$
has degree $p^{n - 1}$ over $k$. We have to show that $a_n$ does not
map to a $p$th power in $k(a_1^{1/p}, \ldots, a_{n - 1}^{1/p})$.
If it does then we can write
\begin{align*}
a_n & =
\left(\sum\nolimits_{I = (i_1, \ldots, i_{n - 1}),\ 0 \leq i_j \leq p - 1}
\lambda_I a_1^{i_1/p} \ldots a_{n - 1}^{i_{n - 1}/p}\right)^p \\
& = \sum\nolimits_{I = (i_1, \ldots, i_{n - 1}),\ 0 \leq i_j \leq p - 1}
\lambda_I^p a_1^{i_1} \ldots a_{n - 1}^{i_{n - 1}}
\end{align*}
Applying $\text{d}$ we see that $\text{d}a_n$ is linearly dependent on
$\text{d}a_i$, $i < n$. This is a contradiction.
\end{proof}


\begin{lemma}
\label{algebra-lemma-derivative-zero-pth-power}
Let $k$ be a perfect field of characteristic $p > 0$.
Let $K/k$ be an extension.
Let $a \in K$. Then $\text{d}a = 0$ in $\Omega_{K/k}$
if and only if $a$ is a $p$th power.
\end{lemma}

\begin{proof}
By Lemma \href{https://stacks.math.columbia.edu/tag/031G}{lemma colimit differentials (Tag 031G)} we see that there exists a subfield
$k \subset L \subset K$ such that $L/k$
is a finitely generated field extension and such that
$\text{d}a$ is zero in $\Omega_{L/k}$.
Hence we may assume that $K$ is a finitely generated field extension
of $k$.

\medskip\noindent
Choose a transcendence basis $x_1, \ldots, x_r \in K$
such that $K$ is finite separable over $k(x_1, \ldots, x_r)$.
This is possible by the definitions, see
Definitions \ref{algebra-definition-perfect} and
\ref{algebra-definition-separable-field-extension}.
We remark that the result holds for the purely transcendental
subfield $k(x_1, \ldots, x_r) \subset K$.
Namely,
$$
\Omega_{k(x_1, \ldots, x_r)/k} =
\bigoplus\nolimits_{i = 1}^r k(x_1, \ldots, x_r) \text{d}x_i
$$
and any rational function all of whose partial derivatives are zero
is a $p$th power. Moreover, we also have
$$
\Omega_{K/k} =
\bigoplus\nolimits_{i = 1}^r K\text{d}x_i
$$
since $k(x_1, \ldots, x_r) \subset K$ is finite separable
(computation omitted). Suppose $a \in K$ is an element such that
$\text{d}a = 0$ in the module of differentials. By our choice of $x_i$ we
see that the minimal polynomial $P(T) \in k(x_1, \ldots, x_r)[T]$
of $a$ is separable. Write
$$
P(T) = T^d + \sum\nolimits_{i = 1}^d a_i T^{d - i}
$$
and hence
$$
0 = \text{d}P(a) = \sum\nolimits_{i = 1}^d a^{d - i}\text{d}a_i
$$
in $\Omega_{K/k}$. By the description of
$\Omega_{K/k}$ above and the fact that $P$ was the minimal
polynomial of $a$, we see that this implies $\text{d}a_i = 0$.
Hence $a_i = b_i^p$ for each $i$. Therefore by
Fields, Lemma \href{https://stacks.math.columbia.edu/tag/031V}{lemma pth root (Tag 031V)}
we see that $a$ is a $p$th power.
\end{proof}


\begin{lemma}
\label{lemma-transitive-lci-at-end}
Let $A \to B \to C$ be ring maps. Assume $B \to C$ is a local complete
intersection homomorphism. Choose a presentation
$\alpha : A[x_s, s \in S] \to B$ with kernel $I$. Choose a presentation
$\beta : B[y_1, \ldots, y_m] \to C$ with kernel $J$. Let
$\gamma : A[x_s, y_t] \to C$ be the induced presentation of $C$ with kernel
$K$. Then we get a canonical commutative diagram
$$
\xymatrix{
0 \ar[r] &
\Omega_{A[x_s]/A} \otimes C \ar[r] &
\Omega_{A[x_s, y_t]/A} \otimes C \ar[r] &
\Omega_{B[y_t]/B} \otimes C \ar[r] &
0 \\
0 \ar[r] &
I/I^2 \otimes C \ar[r] \ar[u] &
K/K^2 \ar[r] \ar[u] &
J/J^2 \ar[r] \ar[u] &
0
}
$$
with exact rows. In particular, the six term exact sequence of
Algebra, Lemma \href{https://stacks.math.columbia.edu/tag/00S2}{lemma exact sequence NL (Tag 00S2)}
can be completed with a zero on the left, i.e., the sequence
$$
0 \to H_1(\NL_{B/A} \otimes_B C) \to
H_1(L_{C/A}) \to
H_1(L_{C/B}) \to
\Omega_{B/A} \otimes_B C \to
\Omega_{C/A} \to
\Omega_{C/B} \to 0
$$
is exact.
\end{lemma}

\begin{proof}
The only thing to prove is the injectivity of the map
$I/I^2 \otimes C \to K/K^2$.
By assumption the ideal $J$ is Koszul-regular.
Hence we have $IA[x_s, y_j] \cap K^2 = IK$ by
Lemma \ref{lemma-conormal-sequence-H1-regular-ideal}.
This means that the kernel of $K/K^2 \to J/J^2$ is
isomorphic to $IA[x_s, y_j]/IK$. Since
$I/I^2 \otimes_A C = IA[x_s, y_j]/IK$ by right exactness
of tensor product, this provides us with the desired injectivity of
$I/I^2 \otimes_A C \to K/K^2$.
\end{proof}


\begin{lemma}
\label{lemma-transitive-colimit-lci-at-end}
Let $A \to B \to C$ be ring maps.
If $B \to C$ is a filtered colimit of local complete intersection
homomorphisms then the conclusion of
Lemma \ref{lemma-transitive-lci-at-end}
remains valid.
\end{lemma}

\begin{proof}
Follows from
Lemma \ref{lemma-transitive-lci-at-end}
and
Algebra, Lemma \ref{algebra-lemma-colimits-NL}.
\end{proof}


\begin{lemma}
\label{lemma-conormal-sequence-H1-regular-ideal}
Let $A$ be a ring. Let $I \subset J \subset A$ be ideals.
Assume that $J/I \subset A/I$ is a $H_1$-regular ideal.
Then $I \cap J^2 = IJ$.
\end{lemma}

\begin{proof}
Follows immediately from Lemma \ref{lemma-conormal-sequence-H1-regular}
by localizing.
\end{proof}


\begin{lemma}
\label{lemma-conormal-sequence-H1-regular}
Let $A$ be a ring. Let $I \subset J \subset A$ be ideals.
Assume that $J/I \subset A/I$ is generated by an $H_1$-regular sequence.
Then $I \cap J^2 = IJ$.
\end{lemma}

\begin{proof}
To prove this choose $g_1, \ldots, g_m \in J$
whose images in $A/I$ form a $H_1$-regular sequence which generates $J/I$.
In particular $J = I + (g_1, \ldots, g_m)$.
Suppose that $x \in I \cap J^2$. Because $x \in J^2$ can write
$$
x =
\sum a_{ij} g_ig_j +
\sum a_j g_j +
a
$$
with $a_{ij} \in A$, $a_j \in I$ and $a \in I^2$.
Then $\sum a_{ij}g_ig_j \in I \cap (g_1, \ldots, g_m)$
hence by
Lemma \ref{lemma-H1-regular-in-quotient}
we see that $\sum a_{ij}g_ig_j \in I(g_1, \ldots, g_m)$.
Thus $x \in IJ$ as desired.
\end{proof}


\begin{lemma}
\label{lemma-H1-regular-in-quotient}
Let $A$ be a ring. Let $I \subset A$ be an ideal.
Let $g_1, \ldots, g_m$ be a sequence in $A$ whose image in
$A/I$ is $H_1$-regular. Then $I \cap (g_1, \ldots, g_m) =
I(g_1, \ldots, g_m)$.
\end{lemma}

\begin{proof}
Consider the exact sequence of complexes
$$
0 \to I \otimes_A K_\bullet(A, g_1, \ldots, g_m)
\to K_\bullet(A, g_1, \ldots, g_m) \to
K_\bullet(A/I, g_1, \ldots, g_m) \to 0
$$
Since the complex on the right has $H_1 = 0$ by assumption we
see that
$$
\Coker(I^{\oplus m} \to I)
\longrightarrow
\Coker(A^{\oplus m} \to A)
$$
is injective. This is equivalent to the assertion of the lemma.
\end{proof}


\begin{lemma}
\label{algebra-lemma-localize-NL}
Let $A \to B$ be a ring map. Let $S \subset B$ be a multiplicative subset.
The canonical map $\NL_{B/A} \otimes_B S^{-1}B \to \NL_{S^{-1}B/A}$
is a quasi-isomorphism.
\end{lemma}

\begin{proof}
We have $S^{-1}B = \colim_{g \in S} B_g$ where we think of $S$
as a directed set (ordering by divisibility), see
Lemma \href{https://stacks.math.columbia.edu/tag/00CR}{lemma localization colimit (Tag 00CR)}.
By Lemma \href{https://stacks.math.columbia.edu/tag/08JZ}{lemma principal localization NL (Tag 08JZ)} each of the maps
$\NL_{B/A} \otimes_B B_g \to \NL_{B_g/A}$
are quasi-isomorphisms.
The lemma follows from Lemma \ref{algebra-lemma-colimits-NL}.
\end{proof}


\begin{lemma}
\label{algebra-lemma-colimits-NL}
Let $R_\lambda \to S_\lambda$ be a system of ring maps over
the directed set $\Lambda$.
Set $R = \colim R_\lambda$ and $S = \colim S_\lambda$.
Then $\NL_{S/R} = \colim \NL_{S_\lambda/R_\lambda}$.
\end{lemma}

\begin{proof}
Recall that $\NL_{S/R}$ is the complex
$I/I^2 \to \bigoplus_{s \in S} S\text{d}[s]$ where $I \subset R[S]$
is the kernel of the canonical presentation $R[S] \to S$.
Now it is clear that $R[S] = \colim R_\lambda[S_\lambda]$ and similarly
that $I = \colim I_\lambda$ where
$I_\lambda = \Ker(R_\lambda[S_\lambda] \to S_\lambda)$.
Hence the lemma is clear.
\end{proof}


\begin{lemma}
\label{algebra-lemma-geometrically-regular}
Let $k$ be a field. Let $A$ be a $k$-algebra.
Assume $A$ is Noetherian.
The following properties of $A$ are equivalent:
\begin{enumerate}
\item $k' \otimes_k A$ is regular for every finitely generated field
extension $k'/k$, and
\item $k' \otimes_k A$ is regular for every finite purely inseparable
extension $k'/k$.
\end{enumerate}
Here regular ring is as in Definition \ref{algebra-definition-regular}.
\end{lemma}

\begin{proof}
The lemma makes sense by the remarks preceding the lemma.
It is clear that (1) $\Rightarrow$ (2).

\medskip\noindent
Assume (2) and let $K/k$ be a finitely generated field extension.
By Lemma \href{https://stacks.math.columbia.edu/tag/030R}{lemma make separable (Tag 030R)} we can find a diagram
$$
\xymatrix{
K \ar[r] & K' \\
k \ar[u] \ar[r] & k' \ar[u]
}
$$
where $k'/k$, $K'/K$ are finite purely inseparable field
extensions such that $K'/k'$ is separable. By
Lemma \href{https://stacks.math.columbia.edu/tag/037X}{lemma localization smooth separable (Tag 037X)}
there exists a smooth $k'$-algebra $B$ such that $K'$ is the
fraction field of $B$. Now we can argue as follows:
Step 1: $k' \otimes_k A$ is a regular ring because we assumed (2).
Step 2: $B \otimes_{k'} k' \otimes_k A$ is a regular ring as
$k' \otimes_k A \to B \otimes_{k'} k' \otimes_k A$ is smooth
(Lemma \href{https://stacks.math.columbia.edu/tag/00T4}{lemma base change smooth (Tag 00T4)})
and ascent of regularity along smooth maps
(Lemma \href{https://stacks.math.columbia.edu/tag/07NF}{lemma regular goes up (Tag 07NF)}).
Step 3. $K' \otimes_{k'} k' \otimes_k A = K' \otimes_k A$ is
a regular ring as it is a localization of a regular ring
(immediate from the definition).
Step 4. Finally $K \otimes_k A$ is a regular ring by descent of
regularity along the faithfully flat ring map
$K \otimes_k A \to K' \otimes_k A$ (Lemma \href{https://stacks.math.columbia.edu/tag/07NG}{lemma descent regular (Tag 07NG)}).
This proves the lemma.
\end{proof}


\noindent
Let $R \to S$ be a ring map. Denote $R[S]$ the polynomial ring
whose variables are the elements $s \in S$. Let's denote $[s] \in R[S]$
the variable corresponding to $s \in S$. Thus $R[S]$ is a free
$R$-module on the basis elements $[s_1] \ldots [s_n]$ where
$s_1, \ldots, s_n$ ranges over all unordered sequences of elements of $S$.
There is a canonical surjection
\begin{equation}
\label{algebra-equation-canonical-presentation}
R[S] \longrightarrow S,\quad [s] \longmapsto s
\end{equation}
whose kernel we denote $I \subset R[S]$. It is a simple observation that
$I$ is generated by the elements
$[s + s'] - [s] - [s']$, $[s][s'] - [ss']$ and $[r] - r$.
According to Lemma \href{https://stacks.math.columbia.edu/tag/00RU}{lemma differential seq (Tag 00RU)}
there is a canonical map
\begin{equation}
\label{algebra-equation-naive-cotangent-complex}
I/I^2 \longrightarrow \Omega_{R[S]/R} \otimes_{R[S]} S
\end{equation}
whose cokernel is canonically isomorphic to $\Omega_{S/R}$. Observe that
the $S$-module $\Omega_{R[S]/R} \otimes_{R[S]} S$ is free on the generators
$\text{d}[s]$.


\noindent
Let $k$ be a field. Let $(A, \mathfrak m, K)$ be a Noetherian local
$k$-algebra. The Jacobi-Zariski sequence
(Algebra, Lemma \href{https://stacks.math.columbia.edu/tag/00S2}{lemma exact sequence NL (Tag 00S2)})
is a canonical exact sequence
$$
H_1(L_{K/k}) \to \mathfrak m/\mathfrak m^2 \to \Omega_{A/k} \otimes_A K \to
\Omega_{K/k} \to 0
$$
because $H_1(L_{K/A}) = \mathfrak m/\mathfrak m^2$ by
Algebra, Lemma \href{https://stacks.math.columbia.edu/tag/07BP}{lemma NL surjection (Tag 07BP)}.
We will show that exactness on the left of this sequence characterizes
whether or not a regular local ring $A$ is geometrically regular over $k$.
We will link this to the notion of formal smoothness in
Section \href{https://stacks.math.columbia.edu/tag/07EK}{section regular fs (Tag 07EK)}.


\begin{proposition}
\label{proposition-characterization-geometrically-regular}
Let $k$ be a field of characteristic $p > 0$.
Let $(A, \mathfrak m, K)$ be a Noetherian local
$k$-algebra. The following are equivalent
\begin{enumerate}
\item $A$ is geometrically regular over $k$,
\item for all $k \subset k' \subset k^{1/p}$
finite over $k$ the ring $A \otimes_k k'$ is regular,
\item $A$ is regular and the canonical map
$H_1(L_{K/k}) \to \mathfrak m/\mathfrak m^2$ is injective, and
\item $A$ is regular and the map
$\Omega_{k/\mathbf{F}_p} \otimes_k K \to \Omega_{A/\mathbf{F}_p} \otimes_A K$
is injective.
\end{enumerate}
\end{proposition}

\begin{proof}
Proof of (3) $\Rightarrow$ (1).
Assume (3). Let $k'/k$ be a finite purely inseparable extension.
Set $A' = A \otimes_k k'$. This is a local ring with maximal ideal
$\mathfrak m'$. Set $K' = A'/\mathfrak m'$. We get a commutative
diagram
$$
\xymatrix{
0 \ar[r] &
H_1(L_{K/k}) \otimes K' \ar[r] \ar[d]_\beta &
\mathfrak m/\mathfrak m^2 \otimes K' \ar[r] \ar[d] &
\Omega_{A/k} \otimes_A K' \ar[r] \ar[d]_{\cong} &
\Omega_{K/k} \otimes K' \ar[r] \ar[d]_\alpha &
0
\\
 &
H_1(L_{K'/k'}) \ar[r] &
\mathfrak m'/(\mathfrak m')^2 \ar[r] &
\Omega_{A'/k'} \otimes_{A'} K' \ar[r] &
\Omega_{K'/k'} \ar[r] &
0
}
$$
with exact rows. The third vertical arrow is an isomorphism by base
change for modules of differentials
(Algebra, Lemma \href{https://stacks.math.columbia.edu/tag/00RV}{lemma differentials base change (Tag 00RV)}).
Thus $\alpha$ is surjective. By
Lemma \ref{lemma-gamma-commutative-diagram} we have
$$
\dim \Ker(\alpha) - \dim \Ker(\beta) + \dim \Coker(\beta) = 0
$$
(and these dimensions are all finite). A diagram chase shows that
$\dim \mathfrak m'/(\mathfrak m')^2 \leq \dim \mathfrak m/\mathfrak m^2$.
However, since $A \to A'$ is finite flat we see that
$\dim(A) = \dim(A')$, see
Algebra, Lemma \href{https://stacks.math.columbia.edu/tag/00OM}{lemma dimension base fibre total (Tag 00OM)}.
Hence $A'$ is regular by definition.

\medskip\noindent
Equivalence of (3) and (4). Consider the Jacobi-Zariski sequences
for rows of the commutative diagram
$$
\xymatrix{
\mathbf{F}_p \ar[r] & A \ar[r] & K \\
\mathbf{F}_p \ar[r] \ar[u] & k \ar[r] \ar[u] & K \ar[u]
}
$$
to get a commutative diagram
$$
\xymatrix{
0 \ar[r] &
\mathfrak m/\mathfrak m^2 \ar[r] &
\Omega_{A/\mathbf{F}_p} \otimes_A K \ar[r] &
\Omega_{K/\mathbf{F}_p} \ar[r] & 0 & \\
0 \ar[r] &
H_1(L_{K/k}) \ar[r] \ar[u] &
\Omega_{k/\mathbf{F}_p} \otimes_k K \ar[r] \ar[u] &
\Omega_{K/\mathbf{F}_p} \ar[r] \ar[u] &
\Omega_{K/k} \ar[r] \ar[u] &
0
}
$$
with exact rows. We have used that $H_1(L_{K/A}) = \mathfrak m/\mathfrak m^2$
and that $H_1(L_{K/\mathbf{F}_p}) = 0$ as $K/\mathbf{F}_p$ is separable, see
Algebra, Proposition
\href{https://stacks.math.columbia.edu/tag/0322}{proposition characterize separable field extensions (Tag 0322)}.
Thus it is clear that the kernels of
$H_1(L_{K/k}) \to \mathfrak m/\mathfrak m^2$
and
$\Omega_{k/\mathbf{F}_p} \otimes_k K \to \Omega_{A/\mathbf{F}_p} \otimes_A K$
have the same dimension.

\medskip\noindent
Proof of (2) $\Rightarrow$ (4) following Faltings, see
\href{https://stacks.math.columbia.edu/bibliography}{Bibliography: Faltings-einfacher}. Let $a_1, \ldots, a_n \in k$ be elements
such that $\text{d}a_1, \ldots, \text{d}a_n$ are linearly independent
in $\Omega_{k/\mathbf{F}_p}$. Consider the field extension
$k' = k(a_1^{1/p}, \ldots, a_n^{1/p})$. By
Algebra, Lemma \ref{algebra-lemma-size-extension-pth-roots}
we see that $k' = k[x_1, \ldots, x_n]/(x_1^p - a_1, \ldots, x_n^p - a_n)$.
In particular we see that the naive cotangent complex of $k'/k$
is homotopic to the complex
$\bigoplus_{j = 1, \ldots, n} k' \rightarrow \bigoplus_{i = 1, \ldots, n} k'$
with the zero differential as
$\text{d}(x_j^p - a_j) = 0$ in $\Omega_{k[x_1, \ldots, x_n]/k}$.
Set $A' = A \otimes_k k'$ and $K' = A'/\mathfrak m'$ as above.
By Algebra, Lemma \href{https://stacks.math.columbia.edu/tag/00S4}{lemma change base NL (Tag 00S4)}
we see that $\NL_{A'/A}$ is homotopy equivalent to the complex
$\bigoplus_{j = 1, \ldots, n} A' \rightarrow \bigoplus_{i = 1, \ldots, n} A'$
with the zero differential, i.e., $H_1(L_{A'/A})$ and
$\Omega_{A'/A}$ are free of rank $n$. The Jacobi-Zariski sequence for
$\mathbf{F}_p \to A \to A'$ is
$$
H_1(L_{A'/A}) \to \Omega_{A/\mathbf{F}_p} \otimes_A A'
\to \Omega_{A'/\mathbf{F}_p} \to \Omega_{A'/A} \to 0
$$
Using the presentation $A[x_1, \ldots, x_n] \to A'$ with
kernel $(x_j^p - a_j)$ we see, unwinding the maps in
Algebra, Lemma \href{https://stacks.math.columbia.edu/tag/00S2}{lemma exact sequence NL (Tag 00S2)},
that the $j$th basis vector of $H_1(L_{A'/A})$ maps to
$\text{d}a_j \otimes 1$ in $\Omega_{A/\mathbf{F}_p} \otimes A'$.
As $\Omega_{A'/A}$ is free (hence flat) we get on tensoring with $K'$
an exact sequence
$$
K'^{\oplus n} \to \Omega_{A/\mathbf{F}_p} \otimes_A K'
\xrightarrow{\beta} \Omega_{A'/\mathbf{F}_p} \otimes_{A'} K' \to
K'^{\oplus n} \to 0
$$
We conclude that the elements $\text{d}a_j \otimes 1$ generate
$\Ker(\beta)$ and we have to show that are linearly independent, i.e.,
we have to show $\dim(\Ker(\beta)) = n$.
Consider the following big diagram
$$
\xymatrix{
0 \ar[r] &
\mathfrak m'/(\mathfrak m')^2 \ar[r] &
\Omega_{A'/\mathbf{F}_p} \otimes K' \ar[r] &
\Omega_{K'/\mathbf{F}_p} \ar[r] & 0 \\
0 \ar[r] &
\mathfrak m/\mathfrak m^2 \otimes K' \ar[r] \ar[u]^\alpha &
\Omega_{A/\mathbf{F}_p} \otimes K' \ar[r] \ar[u]^\beta &
\Omega_{K/\mathbf{F}_p} \otimes K' \ar[r] \ar[u]^\gamma & 0
}
$$
By Lemma \ref{lemma-cartier-equality} and the Jacobi-Zariski sequence for
$\mathbf{F}_p \to K \to K'$ we see that the kernel and cokernel of
$\gamma$ have the same finite dimension.
By assumption $A'$ is regular (and of the same dimension as $A$, see
above) hence the kernel and cokernel of $\alpha$ have the same dimension.
It follows that the kernel and cokernel of $\beta$ have the same
dimension which is what we wanted to show.

\medskip\noindent
The implication (1) $\Rightarrow$ (2) is trivial. This finishes
the proof of the proposition.
\end{proof}
\end{document}
